\begin{papercard}{cory2026}
\cf{Stages} \stagetag{L} (word-level LLM annotation pipeline with passage classifier, retrieval and self-correction; seven models compared), \stagetag{C} (21 GDPR transparency requirements from Art.\ 13 and 14 as annotation scheme), \stagetag{D} (crawl of Android app policies for apps available in Germany). Accepted to PoPETs 2026 (TU Berlin); arXiv v2 used here. Added from the second search round of 23 October 2026: the closest published design to the thesis's span-level, GDPR-framed coding.

\cf{Aim and pipeline} Annotate the exact words of a policy that address each GDPR transparency requirement. The pipeline combines LLM annotation, a passage-level classifier, retrieval and a self-correction step \pc{cory2026}{1}{Our ap- proach integrates LLM-driven annotation with passage-level clas- sification, retrieval-augmented generation, and a self-correction mechanism to deliver scalable, context-aware annotations across 21 GDPR-derived transparency requirements}. The reasoning for passage-level input: long inputs dilute word-level precision \pc{cory2026}{4}{As input length increases, semantic precision may suffer, reducing annotation accuracy}. A second LLM reviews and corrects the first one's annotations \pc{cory2026}{6}{The self-correction layer refines the preliminary annotations generated by the annotation layer by sys- tematically reviewing the set of annotations}, and guardrails replace invented labels or spans \pc{cory2026}{7}{we implement guardrails within the annotation and self-correction layers that prevent them from generating annota- tions that pertain to made-up labels or text spans}.

\cf{Corpus and ground truth} 703,791 English Android-app policies, narrowed to apps available in Germany \pc{cory2026}{8}{we further refine our popula- tion to include only privacy policies of applications available for download in Germany}; 200 policies form the evaluation set. The ground truth was made by six experts \emph{reviewing LLM-generated annotations}, two per policy \pc{cory2026}{9}{we employed a team of six domain experts from two universities in Germany and the Nether- lands to review the LLM-generated annotations for the sample of 200 privacy policies}; their agreement was moderate \pc{cory2026}{10}{This initial round of independent reviews resulted in a mod- erate inter-reviewer agreement (Krippendorff’s Alpha) of 0.669}, and a three-expert jury settled disagreements, giving 25,260 curated annotations.

\cf{Evaluation and results} Two tiers, passage classification and span quality, over seven models (GPT-4o, GPT-4o mini, Llama-3.3, Qwen-2.5, Mixtral, Phi-4, Gemma-2, DeepSeek-R1 are named in the comparison table) and on OPP-115 \pc{cory2026}{1}{We propose a two-tiered evaluation methodology capturing both passage-level classification and span-level annotation quality}. Decomposition helped most \pc{cory2026}{1}{decomposing the annotation task and integrating targeted retrieval and classification compo- nents significantly improve annotation accuracy, particularly for well-structured requirements}. Short, well-defined labels worked; long, context-dependent or legally ambiguous ones did not, and related categories were confused \pc{cory2026}{13}{LLM-based word-level annota- tion is most effective for short, well-structured requirements such as Controller Name and Data Categories, but struggles with labels involving longer, context-dependent spans or ambiguous legal dis- closures}.

\cf{Limits (card)} The reference was built by correcting model output, so it leans towards what the model proposed (anchoring); omissions the model never suggested are less likely to be caught than in blind coding. English policies only, although the population is apps available in Germany. Transparency requirements are what a policy must state; harms and stance are not covered.

\cf{Use for this thesis} (1) Design evidence for D9: classify passages first, then annotate spans within them; retrieve definitions and examples per label; let a second pass check the first; restrict outputs to valid labels and to text that exists in the source. (2) Expected difficulty: harm, stance and protection spans are ``longer, context-dependent'' labels, the kind this pipeline handled worst, which supports manual review of all retained mentions. (3) Even experts agreed only moderately ($\alpha$ 0.669) on GDPR spans, consistent with Arora et al.\ \pp{arora2022}. (4) Contrast for the validation design: the thesis codes blind before seeing model output, which avoids the anchoring built into a review-of-model-output reference. (5) The 21 requirements (controller, purposes, legal basis, recipients, retention, rights, complaint) can seed Axis A's ``additional details''.
\end{papercard}
